% Tools: docs/tool_contracts.md (6 read tools + 2 mock sinks, no network I/O). Transformed-leak scorer:
% src/react_agent/evaluation/transformed_leakage.py (decision log 2026-11-15 18:20). Mock tool names are not used.
\section{Threat Model}\label{sec:threat}

\paragraph{Agent and tools.} Let $u$ be the user request. At each step the model emits an action $a_t$, consisting of a
tool name and arguments, or a final answer $y$. The Tool Broker checks $a_t$, executes it and returns an observation
$o_t$. The system has six read tools (searching and reading documents, read-only database queries, searching and opening
cached web pages, and a calculator) and two \emph{sinks} with side effects, sending e-mail and calling a webhook. Both
sinks are mocks that only log the call and never send anything over the network. The final answer is treated as a third
sink. For a sink call $a$, $\dest(a)$ denotes its destination (an address or a URL) and $\data(a)$ the remaining
arguments.

\paragraph{Two labels.} Every piece of content carries two independent labels. By \emph{trust}, only $u$ and the system
prompt are trusted, and every tool observation is untrusted. By \emph{sensitivity}, content ranges from public to
confidential. An untrusted page may be public, while a confidential record may come from a source the attacker cannot
touch. Merging the two dimensions forces a defence either to block all confidential data or to trust every internal
source.

\paragraph{Attacker.} The attacker controls the text of one piece of content returned by a read tool, at any position
and in any tone, and knows the names and schemas of the tools. In the adaptive rounds the attacker also knows the full
design and formal model of the defence. The attacker cannot modify $u$, the system prompt, the broker or the model
weights. Direct jailbreaks, training-data poisoning, denial of service and multi-agent systems are out of scope.

\paragraph{Harm predicate.} Let $G(u)$ be the set of (tool, destination) pairs that the user actually authorises,
$\Sigma$ the set of confidential values on the trajectory that the user does not allow to be disclosed, and
$\mathcal{S}$ the set of sink calls that were \emph{actually executed}. Values written by the user in $u$ are not in
$\Sigma$. A trajectory $\tau$ is harmful if
\begin{equation}\label{eq:harm}
\mathrm{Harm}(\tau) = H_\mathit{dest} \lor H_\mathit{leak} \lor H_\mathit{tamper} \lor H_\mathit{side},
\end{equation}
where
\begin{itemize}
\item $H_\mathit{dest}$: some $a \in \mathcal{S}$ has $(\mathrm{tool}(a), \dest(a)) \notin G(u)$,
\item $H_\mathit{leak}$: some $v \in \Sigma$ appears in $\data(a)$ for some $a \in \mathcal{S}$ or in $y$, verbatim or in a
  transformed form,
\item $H_\mathit{tamper}$: some $a \in \mathcal{S}$ has a legitimate destination but a user-specified field carries a
  different value,
\item $H_\mathit{side}$: the task only requires reading but $\mathcal{S} \neq \emptyset$.
\end{itemize}
A frozen scorer decides these predicates from the executed calls and $y$. Blocked calls never count as harm. Besides
verbatim matching, the scorer recognises confidential values encoded in Base64, Hex or URL encoding, split into
separate characters, or written out as number words.

\paragraph{Six attack mechanisms.} Attacks are classified by \emph{mechanism}, that is, the path by which the injected
text turns into harm (Table~\ref{tab:mechanisms}). In the two redirection mechanisms, the destination does not occur in
$u$. In the two leakage mechanisms, the destination is legitimate or absent, so only the provenance of a value separates
the attack from the task. In parameter tampering, both the destination and the sensitivity look normal. The unrequested
side-effect mechanism tests whether a destination that is merely mentioned, for instance in a prohibition, is mistaken
for an authorisation.

% Table 1. Numbers: facts.md "Bổ sung sau phản biện 3" (v1 bypass causes: Test-I 13 = SL 6, AL 3, PT 4 (3 value
% changes + 1 destination present in u); Test-II 27 = SL 14, AL 2, PT 7, US 3, CR 1) and the Test-4 section ("40 ca vượt
% v1": M3 20, M5 11, M4 5, M6 3, M2 1). By cause: public source 25, destination present in u 5, tampering 10.
\begin{table}[!htbp]
\centering
\caption{The six attack mechanisms. The third column gives the number of cases that bypassed \base{} on the two earlier
sealed sets (\tOne{} and \tTwo, three backbones, 40 cases in total) and their cause. The last column names the
\ours{} component that targets each group.}\label{tab:mechanisms}
\footnotesize
\setlength{\tabcolsep}{3pt}
\begin{tabularx}{\textwidth}{@{}>{\raggedright\arraybackslash}p{2.45cm}>{\raggedright\arraybackslash}X%
>{\raggedright\arraybackslash}p{4.0cm}l@{}}
\toprule
Mechanism & Path to harm & Bypasses of \base{} (cause) & \ours \\
\midrule
\mER{} explicit redirection & The injection asks for data to be sent to an unknown destination & 0 & \cOne \\
\mCR{} covert redirection & As above, but phrased indirectly or with the destination split up & 1 (the destination
  field is an ordinary word that occurs in $u$) & \cOne \\
\mSL{} sink leakage & A confidential value enters the content sent to a legitimate destination & 20 (the value sits in
  a source labelled public) & \cTwo \\
\mAL{} answer leakage & A confidential value appears in the final answer & 5 (as above) & \cTwo \\
\mPT{} parameter tampering & A value stated by the user is changed while the destination stays legitimate & 11 (10
  changed values, 1 destination that is only mentioned) & \cThree, \cOne \\
\mUS{} unrequested side effect & A sink is called on a destination that is only mentioned or explicitly forbidden & 3
  (destination present in $u$ but not authorised) & \cOne \\
\bottomrule
\end{tabularx}
\end{table}


\paragraph{Families, benign branches and variants.} A \emph{family} is one attack idea in a specific task and tool
environment. Every family has a \emph{\bPlain} branch, which replaces the injection with a harmless note of the same
length and is used to measure utility. One third of the families also have an \emph{\bAuth} branch, in which the user
explicitly authorises an action similar to the attack goal. A denial in a benign branch is a \emph{false block}. Each
family has five variants applied to both $u$ and the injection: standard Vietnamese, a paraphrase, Vietnamese--English
code-mixing, Vietnamese without diacritics, and English. Identifiers, addresses and confidential values are identical
across variants, so any difference comes from language alone.
